% =====================================================================
% Anexo D: analisis arquitectonico y computacional de los componentes.
% =====================================================================

\section*{Anexo D. Análisis arquitectónico y computacional de los componentes}
\label{sec:analisis-componentes}

% Numeracion propia del anexo: figuras, tablas y ecuaciones como D.n
\setcounter{figure}{0}\renewcommand{\thefigure}{D.\arabic{figure}}
\setcounter{table}{0}\renewcommand{\thetable}{D.\arabic{table}}
\setcounter{equation}{0}\renewcommand{\theequation}{D.\arabic{equation}}

Este anexo examina los dos componentes del modelo en términos de su realización:
qué tipo de capa emplea cada uno, cuántos parámetros concentra, cuánta memoria
exige, en qué régimen opera y hasta dónde admite distribución. Se apoya en la
inspección directa del código liberado por ambos grupos. Sus tres resultados se
recogen en la Sección~\ref{sec:analisis-resumen} del marco teórico.

\subsection*{D.1 Revisión de las implementaciones de referencia}
\label{sec:rev-implementaciones}

\paragraph*{Sonata} El repositorio se distribuye como material de inferencia y
visualización, no de preentrenamiento. La arquitectura instanciada reproduce la
configuración descrita, cinco etapas de profundidades $(3,3,3,12,3)$ y anchuras
$(48,96,192,384,512)$, y el recuento analítico de parámetros coincide con el
total que el modelo declara al cargarse (Tabla~\ref{tab:params-sonata}). La
dependencia de FlashAttention es opcional y el código provee una ruta
alternativa explícita.

\paragraph*{HyperNCD} El repositorio contiene dos realizaciones del razonamiento
por hipergrafo. La primera reproduce la formulación del artículo, con matriz de
incidencia, convolución de hipergrafo y actualización dinámica de hiperaristas,
pero su inspección revela que no es invocada desde ninguna ruta de ejecución,
que las matrices de peso de sus capas se instancian dentro de la propia función
en cada llamada y por tanto no reciben gradiente, y que el procedimiento de
actualización dinámica reproduce término por término el de construcción inicial.
La segunda, definida dentro del paso de entrenamiento, es la efectivamente
ejecutada: opera sobre regiones, combina las similitudes de las
Ecs.~\eqref{eq:coseno} y~\eqref{eq:gauss} según la Ec.~\eqref{eq:afinidad},
retiene $M=8$ vecinos y aplica una única agregación normalizada, sin pesos
entrenables ni convolución de hipergrafo en el sentido del artículo.

De ahí la decisión que el trabajo adopta: la línea base es la implementación que
se ejecuta y la sustitución de la Ec.~\eqref{eq:hibrido} se aplica sobre su
descriptor geométrico de seis dimensiones.

\subsection*{D.2 Arquitectura de Sonata}
\label{sec:arq-sonata}

Antes de integrar el \emph{encoder} conviene caracterizarlo, porque de su
estructura se derivan tanto el costo computacional del método como el margen de
optimización disponible. Sonata es un Point Transformer~V3 (PTv3) en
configuración \emph{encoder-only}: cinco etapas jerárquicas, sin decodificador,
con $108{,}46$~M de parámetros. La ausencia de decodificador es deliberada, ya
que el objetivo no es predecir etiquetas sino producir una representación
$f_{\mathrm{ssl}}$, y obliga a reconstruir la escala original mediante la
propagación inversa del submuestreo.

\begin{figure}[H]
\centering
\begin{tikzpicture}[
  font=\scriptsize,
  et/.style={draw, rounded corners, align=center, fill=orange!20,
             minimum height=1.05cm},
  io/.style={draw, rounded corners, align=center, fill=black!6,
             minimum width=8.6cm, minimum height=0.55cm},
  pre/.style={draw, rounded corners, align=center, fill=green!16,
              minimum width=8.6cm, minimum height=0.55cm},
  itf/.style={font=\tiny, text=black!65, fill=white, inner sep=1.5pt},
  anot/.style={font=\tiny, text=black!60},
  fl/.style={-{Latex[length=1.6mm]}, thick},
  tap/.style={-{Latex[length=1.4mm]}, dashed, black!55}
]
\node[io] (in) {\textbf{Interfaz de entrada:} $N$ puntos $\times$ 9 canales (coordenadas, color, normal)};
\node[pre, below=2.6mm of in] (pre)
  {\textbf{Preprocesamiento:} serialización (orden Z / Hilbert) $+$ \emph{embedding} lineal $9\rightarrow48$};

\node[et, minimum width=8.0cm, below=8mm of pre] (e1)
  {\textbf{Etapa 1} $\cdot$ 3 bloques\\48 canales $\cdot$ 3 cabezas $\cdot$ resolución $N$\\{\tiny\textcolor{black!60}{0{,}3\,\% de los parámetros}}};
\node[et, minimum width=7.0cm, below=9mm of e1] (e2)
  {\textbf{Etapa 2} $\cdot$ 3 bloques\\96 canales $\cdot$ 6 cabezas $\cdot$ resolución $N/2$\\{\tiny\textcolor{black!60}{1{,}0\,\% de los parámetros}}};
\node[et, minimum width=6.2cm, below=9mm of e2] (e3)
  {\textbf{Etapa 3} $\cdot$ 3 bloques\\192 canales $\cdot$ 12 cabezas $\cdot$ resolución $N/4$\\{\tiny\textcolor{black!60}{4{,}1\,\% de los parámetros}}};
\node[et, minimum width=5.4cm, below=9mm of e3, fill=orange!36] (e4)
  {\textbf{Etapa 4} $\cdot$ 12 bloques\\384 canales $\cdot$ 24 cabezas $\cdot$ resolución $N/8$\\{\tiny\textcolor{black!70}{\textbf{65{,}3\,\% de los parámetros}}}};
\node[et, minimum width=4.6cm, below=9mm of e4, fill=orange!28] (e5)
  {\textbf{Etapa 5} $\cdot$ 3 bloques\\512 canales $\cdot$ 32 cabezas $\cdot$ resol.\ $N/16$\\{\tiny\textcolor{black!70}{29{,}0\,\% de los parámetros}}};

\draw[fl] (pre) -- (e1);
\draw[fl] (e1) -- node[itf] {\emph{pooling} $\times\tfrac{1}{2}$ $\cdot$ lineal $48\rightarrow96$} (e2);
\draw[fl] (e2) -- node[itf] {\emph{pooling} $\times\tfrac{1}{2}$ $\cdot$ lineal $96\rightarrow192$} (e3);
\draw[fl] (e3) -- node[itf] {\emph{pooling} $\times\tfrac{1}{2}$ $\cdot$ lineal $192\rightarrow384$} (e4);
\draw[fl] (e4) -- node[itf] {\emph{pooling} $\times\tfrac{1}{2}$ $\cdot$ lineal $384\rightarrow512$} (e5);

\node[io, below=8mm of e5] (out)
  {\textbf{Interfaz de salida:} $N/16$ puntos $\times$ 512 canales};
\node[io, fill=purple!12, below=2.6mm of out] (fssl)
  {$f_{\mathrm{ssl}}$: \emph{up-cast} de las etapas 3, 4 y 5, $192+384+512=1088$};
\draw[fl] (e5) -- (out);

\coordinate (bus) at (5.25,0);
\foreach \e in {e3, e4, e5}
  \draw[dashed, black!55] (\e.east) -- (\e.east -| bus);
\draw[dashed, black!55] (e3.east -| bus) -- (fssl.east -| bus);
\draw[tap] (fssl.east -| bus) -- (fssl.east);

\node[anot, rotate=90, anchor=south] at ($(e1.west)!0.5!(e5.west)+(-4.55,0)$)
  {anchura $\uparrow$ \quad resolución $\downarrow$};
\node[anot, align=center, anchor=north] at ($(e1.north west)+(-2.6,0.55)$)
  {cada etapa repite\\el mismo bloque\\(Fig.~\ref{fig:bloque-sonata})};
\draw[tap, black!40] ($(e1.north west)+(-1.1,-0.35)$) -- ($(e1.west)+(0,0.1)$);

\node[below=3mm of fssl, font=\scriptsize] (leg)
  {\tikz\node[draw,rounded corners,fill=orange!20,minimum height=2.6mm,minimum width=3.5mm]{};\ etapa \emph{Transformer}\quad
   \tikz\node[draw,rounded corners,fill=green!16,minimum height=2.6mm,minimum width=3.5mm]{};\ preprocesamiento\quad
   \tikz\node[draw,rounded corners,fill=black!6,minimum height=2.6mm,minimum width=3.5mm]{};\ interfaz\quad
   \tikz\node[draw,rounded corners,fill=purple!12,minimum height=2.6mm,minimum width=3.5mm]{};\ salida usada};
\end{tikzpicture}
\caption{Arquitectura del \emph{encoder} de Sonata: cinco etapas jerárquicas sin
decodificador. Cada etapa es un componente del mismo tipo, formado por bloques
idénticos (Fig.~\ref{fig:bloque-sonata}), y queda caracterizada por tres
magnitudes: resolución, anchura y capacidad. Al descender la jerarquía, la
resolución se divide por dos en cada transición mientras la anchura crece de 48
a 512 canales, de modo que el $94{,}3$\,\% de los parámetros se concentra en las
dos últimas etapas. Las etiquetas sobre los conectores declaran la interfaz
entre componentes. Las líneas discontinuas señalan las tres etapas cuyas salidas
se concatenan durante la propagación inversa del submuestreo para formar
$f_{\mathrm{ssl}}$. Elaboración propia.}
\label{fig:arq-sonata}
\end{figure}

La composición de la Figura~\ref{fig:arq-sonata} impone dos propiedades que se
usan más adelante. La primera es que la dependencia entre etapas es estricta:
cada una consume la salida de la anterior, de modo que la profundidad no admite
reparto entre unidades de cómputo y el paralelismo debe buscarse dentro de cada
bloque. La segunda es que no hay convoluciones densas ni capas recurrentes; el
único componente convolucional es la codificación posicional dispersa interna a
cada bloque, cuya estructura detalla la Figura~\ref{fig:bloque-sonata}.

\begin{figure}[H]
\centering
\begin{tikzpicture}[
  font=\scriptsize,
  op/.style={draw, rounded corners, minimum width=4.6cm, minimum height=0.6cm,
             align=center, fill=orange!18},
  cv/.style={draw, rounded corners, minimum width=4.6cm, minimum height=0.6cm,
             align=center, fill=green!16},
  nm/.style={draw, rounded corners, minimum width=4.6cm, minimum height=0.5cm,
             align=center, fill=yellow!22},
  par/.style={draw, rounded corners, minimum width=1.05cm, minimum height=0.75cm,
              align=center, fill=orange!30},
  fl/.style={-{Latex[length=1.8mm]}, thick}
]
\node[cv] (cpe) {Conv.\ dispersa \emph{submanifold} $3\times3\times3$\\\emph{(codificación posicional)}};
\node[nm, below=2.2mm of cpe] (n1) {Normalización};
\node[op, below=2.2mm of n1] (at) {\textbf{Atención serializada multi-cabeza}};
\node[nm, below=2.2mm of at] (n2) {Normalización};
\node[op, below=2.2mm of n2, fill=blue!14] (mlp) {MLP: $C\rightarrow4C\rightarrow C$};
\foreach \a/\b in {cpe/n1, n1/at, at/n2, n2/mlp} \draw[fl] (\a) -- (\b);

\draw[fl, black!55] ($(cpe.north)+(0,0.35)$) -- (cpe);
\draw[fl, black!55] (mlp) -- ++(0,-0.45);
\draw[black!55, rounded corners]
  ($(cpe.north)+(0,0.2)$) -- ++(-2.75,0) |- ($(at.south)+(0,-0.11)$);
\draw[black!55, rounded corners]
  ($(n2.north)+(0,0.14)$) -- ++(2.75,0) |- ($(mlp.south)+(0,-0.11)$);
\node[left=2.9cm of n1, font=\tiny, text=black!55, align=center, rotate=90, anchor=south]
  {residual};

\node[right=1.4cm of at, par] (p1) {$P_1$};
\node[right=1.5mm of p1, par] (p2) {$P_2$};
\node[right=1.5mm of p2, font=\scriptsize] (pd) {$\cdots$};
\node[right=1.5mm of pd, par] (pm) {$P_m$};
\node[above=1.5mm of p2, font=\tiny, align=center, text=black!70, xshift=6mm]
  {$m=N/K$ parches de $K=1024$ puntos\\\textbf{procesados en paralelo}};
\draw[fl, dashed, black!55] (at.east) -- (p1.west);
\node[below=1.5mm of p2, font=\tiny, align=center, text=black!70, xshift=6mm]
  {cada parche: $H$ cabezas en paralelo};
\end{tikzpicture}
\caption{Bloque elemental, repetido 24 veces. Dentro de cada bloque, la atención
se calcula de forma independiente sobre $m=N/K$ parches contiguos de la secuencia
serializada, y dentro de cada parche sobre $H$ cabezas; ambos niveles se resuelven
en paralelo como una única operación densa por lotes. El paralelismo del modelo
reside en esta anchura, no en la profundidad.}
\label{fig:bloque-sonata}
\end{figure}

\paragraph*{Bloque elemental}
Cada uno de los 24 bloques está compuesto por una codificación posicional
condicional implementada como convolución dispersa submanifold de núcleo
$3\times3\times3$, un módulo de atención serializada y un perceptrón multicapa
con razón de expansión 4. Para una anchura $C$, el número de parámetros por
bloque es
\begin{equation}
\begin{aligned}
P(C)&=\underbrace{27C^{2}+C^{2}}_{\text{cod. posicional}}
     +\underbrace{3C^{2}+C^{2}}_{\text{atención}}
     +\underbrace{8C^{2}}_{\text{MLP}}+17C\\[2pt]
    &=40C^{2}+17C .
\end{aligned}
\label{eq:params-bloque}
\end{equation}
La Ec.~\eqref{eq:params-bloque} reproduce exactamente el total de
$108\,461\,328$ parámetros que declara la implementación, lo que valida el
desglose de la Tabla~\ref{tab:params-sonata}.

\begin{table}[H]
\centering
\caption{Distribución de parámetros del \emph{encoder} de Sonata por etapa.}
\label{tab:params-sonata}
\begin{tabular}{@{}lrrrr@{}}
\toprule
\textbf{Etapa} & \textbf{Canales} & \textbf{Bloques} & \textbf{Parámetros} & \textbf{\%} \\
\midrule
1 & 48 & 3 & $278\,928$ & 0{,}3 \\
2 & 96 & 3 & $1\,110\,816$ & 1{,}0 \\
3 & 192 & 3 & $4\,433\,472$ & 4{,}1 \\
4 & 384 & 12 & $70\,857\,216$ & \textbf{65{,}3} \\
5 & 512 & 3 & $31\,483\,392$ & 29{,}0 \\
\emph{Embedding} y submuestreo & & & $297\,504$ & 0{,}3 \\
\midrule
\textbf{Total} & & & $\mathbf{108\,461\,328}$ & 100{,}0 \\
\bottomrule
\end{tabular}
\end{table}

El $94{,}3$\,\% de la capacidad se concentra en las etapas 4 y 5, donde la
resolución espacial es menor y la representación más abstracta. Los pesos ocupan
$0{,}40$~GB en precisión simple, cifra despreciable frente a la memoria de
activaciones que se analiza a continuación.

\paragraph*{Atención serializada y complejidad}
La propiedad que hace viable el modelo sobre nubes de gran tamaño es que la
atención no se calcula sobre vecindarios geométricos. Los puntos se ordenan
mediante una curva de llenado de espacio, de tipo Z o Hilbert, de modo que la
proximidad en $\mathbb{R}^{3}$ se traduce en proximidad dentro de una secuencia
unidimensional; esa secuencia se divide en parches contiguos de $K=1024$ puntos y
la atención se calcula únicamente dentro de cada parche. El costo pasa así de
$O(N^{2})$ a
\begin{equation}
O(N\cdot K),\qquad K=1024\ \text{constante},
\label{eq:complejidad}
\end{equation}
es decir, lineal en el número de puntos. Esta es la diferencia estructural
respecto de arquitecturas basadas en grafos $k$-NN dinámicos, cuyo costo de
construcción de vecindarios crece de forma superlineal y constituye su principal
limitación de escalabilidad.

\paragraph*{Memoria requerida}
Cuando FlashAttention no está disponible, la matriz de atención se materializa de
forma explícita. Su tamaño para una etapa con $N_{s}$ puntos y $H_{s}$ cabezas es
\begin{equation}
M_{s}=N_{s}\cdot H_{s}\cdot K\cdot b ,
\label{eq:memoria}
\end{equation}
con $b$ el tamaño en \emph{bytes} del tipo de dato. Puesto que cada etapa reduce
$N_{s}$ a la mitad mientras duplica $H_{s}$, el producto $N_{s}H_{s}$ permanece
aproximadamente constante y la memoria se distribuye de manera uniforme entre
etapas, como muestra la Tabla~\ref{tab:vram-sonata}.

\begin{table}[H]
\centering
\caption{Estimación de memoria de vídeo por etapa para una sección de
$3\times10^{5}$ puntos tras voxelización, en precisión simple y sin
FlashAttention.}
\label{tab:vram-sonata}
\begin{tabular}{@{}lrrrr@{}}
\toprule
\textbf{Etapa} & \textbf{Puntos} & \textbf{Cabezas} & \textbf{Atención (GB)} & \textbf{QKV (GB)} \\
\midrule
1 & $300\,000$ & 3 & 3{,}43 & 0{,}16 \\
2 & $150\,000$ & 6 & 3{,}43 & 0{,}16 \\
3 & $75\,000$ & 12 & 3{,}43 & 0{,}16 \\
4 & $37\,500$ & 24 & 3{,}43 & 0{,}16 \\
5 & $18\,750$ & 32 & 2{,}29 & 0{,}11 \\
\midrule
Pico por bloque & & & 3{,}59 & \\
Características persistentes & & & 0{,}25 & \\
Pesos del modelo & & & 0{,}40 & \\
\midrule
\textbf{Total estimado} & & & $\mathbf{\approx 4{,}2}$ & \\
\bottomrule
\end{tabular}
\end{table}

La estimación de $4{,}2$~GB se sitúa muy por debajo de los 48~GB disponibles en
la GPU asignada, lo que confirma que la inferencia del \emph{encoder} no es el
factor limitante del proyecto. Se trata de una cota derivada analíticamente a
partir de la Ec.~\eqref{eq:memoria} y del número de puntos supuesto; el
procedimiento de cacheo registra el consumo real mediante
\texttt{max\_memory\_allocated}, de modo que estas cifras serán sustituidas por
mediciones directas.

\paragraph*{Paralelismo y optimización}
Los parches resultantes de la serialización son independientes entre sí, por lo
que se procesan simultáneamente como una dimensión de lote y la carga se resuelve
en operaciones densas de multiplicación de matrices, favorables a la arquitectura
de la GPU. A ese paralelismo interno se añade uno externo: dado que el
\emph{encoder} opera congelado y cada sección de la catedral se procesa de forma
independiente, no existe sincronización de gradientes y el cacheo admite
distribución directa entre nodos mediante arreglos de trabajos del gestor de
colas; la Fig.~\ref{fig:paralelismo} contrasta ambos esquemas de
ejecución.

\begin{figure}[H]
\centering
\begin{tikzpicture}[
  font=\scriptsize,
  sec/.style={draw, minimum height=0.55cm, align=center, fill=orange!22},
  gpu/.style={font=\scriptsize, anchor=east},
  eje/.style={-{Latex[length=1.6mm]}, black!70}
]

%% ===== (a) Ejecucion actual =====
\node[font=\bfseries\scriptsize, anchor=west] at (0,1.35) {(a) Ejecución actual: secuencial};

\node[gpu] at (-0.15,0.55) {GPU 1};
\node[sec, minimum width=1.5cm] (a1) at (0.75,0.55) {$S_1$};
\node[sec, minimum width=1.5cm, right=0mm of a1] (a2) {$S_2$};
\node[sec, minimum width=1.5cm, right=0mm of a2] (a3) {$S_3$};
\node[sec, minimum width=1.5cm, right=0mm of a3] (a4) {$S_4$};

\draw[eje] (0,0.1) -- (7.4,0.1) node[below left=0.5mm and 0mm, font=\tiny] {tiempo};
\draw[black!55] (0,0.2) -- (0,0.0);
\draw[black!55] (6.0,0.2) -- (6.0,0.0);
\node[font=\tiny, below] at (6.0,0.05) {$4t$};
\node[font=\tiny, right=1mm of a4, text=black!70] {3 GPU ociosas};

%% ===== (b) Ejecucion paralela =====
\node[font=\bfseries\scriptsize, anchor=west] at (0,-0.75) {(b) Con arreglo de trabajos: secciones en paralelo};

\node[gpu] at (-0.15,-1.55) {GPU 1};
\node[sec, minimum width=1.5cm] at (0.75,-1.55) {$S_1$};
\node[gpu] at (-0.15,-2.2) {GPU 2};
\node[sec, minimum width=1.5cm] at (0.75,-2.2) {$S_2$};
\node[gpu] at (-0.15,-2.85) {GPU 3};
\node[sec, minimum width=1.5cm] at (0.75,-2.85) {$S_3$};
\node[gpu] at (-0.15,-3.5) {GPU 4};
\node[sec, minimum width=1.5cm] at (0.75,-3.5) {$S_4$};

\draw[eje] (0,-3.95) -- (7.4,-3.95) node[below left=0.5mm and 0mm, font=\tiny] {tiempo};
\draw[black!55] (0,-3.85) -- (0,-4.05);
\draw[black!55] (1.5,-3.85) -- (1.5,-4.05);
\node[font=\tiny, below] at (1.5,-4.0) {$t$};

\draw[{Latex[length=1.6mm]}-{Latex[length=1.6mm]}, black!60]
  (1.6,-2.5) -- (5.9,-2.5)
  node[midway, above, font=\tiny, text=black!70]
  {aceleración $\approx 4\times$ en tiempo de reloj};

\node[align=left, font=\tiny, text=black!70, anchor=north west] at (0,-4.45)
  {El costo por sección no cambia; lo que se reduce es el tiempo total.\\
   El paralelismo es perfecto: el \emph{encoder} está congelado y no hay
   sincronización de gradientes entre secciones.};
\end{tikzpicture}
\caption{Esquema de ejecución del cacheo de características. En (a) las secciones
se procesan uno tras otra sobre una sola unidad; en (b) se distribuyen mediante un
arreglo de trabajos del gestor de colas. La independencia entre secciones,
consecuencia de mantener el \emph{encoder} congelado, hace que la aceleración sea
proporcional al número de unidades asignadas.}
\label{fig:paralelismo}
\end{figure}

Las vías de optimización disponibles, ordenadas por impacto sobre la
Ec.~\eqref{eq:memoria}, son la habilitación de FlashAttention, que evita
materializar la matriz de atención; la reducción del tamaño de parche $K$, que
disminuye la memoria de forma proporcional a costa de contexto por parche; y el
uso de precisión reducida, que actúa sobre el factor $b$. La reducción del tamaño
de vóxel actúa sobre $N_{s}$, pero compromete la detección de elementos finos y
está acotada por los rangos de operación declarados.

\subsection*{D.3 Arquitectura de HyperNCD}
\label{sec:arq-hyperncd}

El segundo componente del modelo híbrido presenta una naturaleza opuesta a la del
\emph{encoder} auto-supervisado: se entrena, opera sobre convoluciones dispersas
en lugar de atención, y su costo se paga en cada época. Esta subsección
caracteriza su estructura a partir de la implementación de referencia, ya que de
ella dependen tanto el punto exacto de inserción del aporte como las
restricciones de paralelización del experimento.

\paragraph*{Red troncal}
La red troncal es una MinkowskiUNet-34C, arquitectura codificador-decodificador
sobre tensores dispersos con conexiones de salto entre niveles simétricos. A
diferencia de PTv3, el mecanismo de agregación es una convolución dispersa de
núcleo $3\times3\times3$ sobre vecindarios geométricos explícitos, no una
atención sobre secuencias serializadas. La configuración declarada en la
implementación es de cuatro niveles descendentes con profundidades
$(2,3,4,6)$ y anchuras $(32,64,128,256)$, y cuatro ascendentes con
profundidades $(2,2,2,2)$ y anchuras $(256,128,96,96)$, sobre una dimensión
inicial de 32 canales.

\begin{figure}[H]
\centering
\begin{tikzpicture}[
  font=\scriptsize,
  cnv/.style={draw, rounded corners, minimum width=2.5cm, minimum height=0.5cm,
              align=center, fill=green!16},
  hd/.style={draw, rounded corners, minimum width=2.9cm, minimum height=0.5cm,
             align=center, fill=blue!14},
  hg/.style={draw, rounded corners, minimum width=3.4cm, minimum height=0.62cm,
             align=center, fill=purple!16},
  dat/.style={draw, rounded corners, minimum width=2.5cm, minimum height=0.45cm,
              align=center, fill=black!6},
  fl/.style={-{Latex[length=1.7mm]}, thick}
]
\node[dat] (in) {Nube voxelizada};
\node[cnv, below=2mm of in] (d1) {Nivel 1: 32 can.};
\node[cnv, below=1.6mm of d1] (d2) {Nivel 2: 64 can.};
\node[cnv, below=1.6mm of d2] (d3) {Nivel 3: 128 can.};
\node[cnv, below=1.6mm of d3, fill=green!30] (d4) {Nivel 4: 256 can.};
\node[cnv, right=1.6cm of d3] (u3) {Nivel 5: 256 can.};
\node[cnv, above=1.6mm of u3] (u2) {Nivel 6: 128 can.};
\node[cnv, above=1.6mm of u2] (u1) {Nivel 7: 96 can.};
\node[cnv, above=1.6mm of u1] (u0) {Nivel 8: 96 can.};
\node[dat, above=2mm of u0] (fs) {$f_{\mathrm{sem}}$ por punto};

\foreach \a/\b in {in/d1, d1/d2, d2/d3, d3/d4} \draw[fl] (\a) -- (\b);
\draw[fl] (d4.east) -- ++(0.55,0) |- (u3.west);
\foreach \a/\b in {u3/u2, u2/u1, u1/u0, u0/fs} \draw[fl] (\a) -- (\b);
\foreach \a/\b in {d3/u3, d2/u2, d1/u1}
  \draw[dashed, black!45] (\a.east) -- (\b.west);
\node[font=\tiny, text=black!55] at ($(d2.east)!0.5!(u2.west)+(0,0.22)$) {saltos};

\node[hg, right=1.5cm of fs, yshift=-0.1cm] (reg) {Agrupamiento en regiones\\(superpuntos)};
\node[hg, below=2.2mm of reg] (geo) {$f_{\mathrm{geo}}$: linealidad, planaridad,\\dispersión y centroide (6\,dim)};
\node[hg, below=2.2mm of geo, fill=purple!30] (hyp) {Hipergrafo: $S_b=\alpha S_g+(1-\alpha)S_s$\\top-$M$, $M=8$, sin parámetros};
\node[hd, below=2.2mm of hyp] (pro) {Cabezas de prototipos};
\node[dat, below=2.2mm of pro, minimum width=3.4cm] (sk) {Sinkhorn-Knopp\\(pseudo-etiquetas)};

\draw[fl] (fs.east) -- (reg.west);
\foreach \a/\b in {reg/geo, geo/hyp, hyp/pro, pro/sk} \draw[fl] (\a) -- (\b);
\draw[fl, red!65, thick] ($(geo.west)+(-0.8,0)$) -- (geo.west)
  node[midway, above, font=\tiny, text=red!70, align=center] {aporte:\\$f_{\mathrm{ssl}}$};
\end{tikzpicture}
\caption{Estructura de HyperNCD. La red troncal dispersa produce
$f_{\mathrm{sem}}$ por punto; el agrupamiento en regiones y el descriptor
geométrico analítico alimentan el hipergrafo, cuya salida se convierte en
pseudo-etiquetas mediante Sinkhorn-Knopp. La flecha roja indica el punto donde
el presente trabajo sustituye el descriptor analítico.}
\label{fig:arq-hyperncd}
\end{figure}

\paragraph*{Razonamiento por hipergrafo}
El hipergrafo no opera sobre puntos individuales sino sobre regiones obtenidas
por agrupamiento previo, lo que reduce el número de nodos en varios órdenes de
magnitud. Para cada par de regiones se calcula la similitud semántica por coseno
de la Ec.~\eqref{eq:coseno} y la geométrica por núcleo gaussiano de la
Ec.~\eqref{eq:gauss}, esta última sobre el descriptor $f_{\mathrm{geo}}$, y
ambas se combinan según la Ec.~\eqref{eq:afinidad}. Las hiperaristas se forman
reteniendo los $M=8$ vecinos de mayor similitud. La
agregación resultante es una única propagación normalizada, sin pesos
entrenables. El descriptor $f_{\mathrm{geo}}$ está compuesto por los tres
índices de forma derivados de los autovalores de la matriz de covarianza local,
es decir linealidad, planaridad y dispersión, concatenados con el centroide de la
región, de modo que su dimensión total es 6. Es precisamente esta representación
de baja dimensionalidad y construcción manual la que la Ec.~\eqref{eq:hibrido}
sustituye por $f_{\mathrm{ssl}}$.

\paragraph*{Costo de entrenamiento}
A diferencia del \emph{encoder} congelado, esta rama requiere almacenar
gradientes y estados del optimizador. Con AdamW, cada parámetro ocupa cuatro
palabras en precisión simple (peso, gradiente y dos momentos), de modo que la
memoria de estado asciende a $16$~bytes por parámetro, frente a los $4$~bytes que
consume el \emph{encoder} en inferencia. A ello se añaden las activaciones
intermedias de la red troncal, que deben conservarse para la retropropagación y
que dominan el consumo total; la implementación mitiga ese costo mediante
submuestreo a $6\times10^{4}$ puntos por nube y lotes de tamaño reducido. La
pérdida total combina el término por puntos y el término por regiones,
\begin{equation}
\mathcal{L}=\mathcal{L}_{\mathrm{cluster}}+\alpha\,\mathcal{L}_{\mathrm{region}} ,
\label{eq:perdida}
\end{equation}
ambos calculados sobre las asignaciones producidas por Sinkhorn-Knopp.

\paragraph*{Restricciones de paralelización}
La configuración de referencia fija una única unidad de cómputo, y tanto la
precisión mixta como el entrenamiento distribuido aparecen desactivados. Esta
limitación no es incidental. El procedimiento de asignación de pseudo-etiquetas
opera por normalización iterativa sobre el lote completo, de modo que fragmentar
el lote entre unidades altera el significado de la asignación salvo que se
introduzca una sincronización explícita del paso de normalización. En
consecuencia, y a diferencia del cacheo de características descrito en la
Fig.~\ref{fig:paralelismo}, el entrenamiento de esta rama no admite distribución
inmediata. El paralelismo disponible se sitúa en un nivel superior: las
particiones experimentales son independientes entre sí y pueden ejecutarse de
forma simultánea, lo que permite aprovechar varias unidades sin modificar el
procedimiento de optimización.

\begin{table}[H]
\centering
\caption{Contraste entre las dos ramas del modelo híbrido.}
\label{tab:contraste-ramas}
\begin{tabular}{@{}p{3.4cm}p{4.4cm}p{4.4cm}@{}}
\toprule
 & \textbf{Sonata (PTv3)} & \textbf{HyperNCD (MinkUNet-34C)} \\
\midrule
Agregación & Atención sobre secuencia serializada & Convolución dispersa $3\times3\times3$ \\
Estructura & Solo codificador, 5 etapas & Codificador-decodificador con saltos \\
Régimen & Congelado, solo inferencia & Entrenado \\
Memoria por parámetro & 4\,bytes & 16\,bytes (AdamW) \\
Frecuencia del costo & Una vez por sección & En cada época \\
Paralelismo entre unidades & Perfecto, sin sincronización & Acoplado por el lote; solo entre particiones \\
\bottomrule
\end{tabular}
\end{table}

Esta asimetría justifica la estrategia de cacheo adoptada: al desplazar el
\emph{encoder} fuera del bucle de optimización, su costo deja de multiplicarse
por el número de épocas y su paralelismo, que es el favorable de los dos, puede
explotarse por completo.

\subsection*{D.4 Escenarios de ejecución}
\label{sec:escenarios}

La viabilidad del método no depende de una única configuración de cómputo, sino
de un conjunto de escenarios cuyo comportamiento conviene anticipar. La
Tabla~\ref{tab:escenarios} los enumera junto con el factor que limita a cada uno.

\begin{table}[H]
\centering
\caption{Escenarios de ejecución considerados para el cacheo de características.}
\label{tab:escenarios}
\begin{tabular}{@{}p{3.5cm}p{2.2cm}p{5.6cm}@{}}
\toprule
\textbf{Escenario} & \textbf{Viable} & \textbf{Factor limitante} \\
\midrule
Portátil, CPU en un hilo & Sí, inviable en la práctica & Las operaciones densas se serializan; sin paralelismo de datos \\
Portátil, CPU multihilo & Marginal & Ancho de banda de memoria, no cómputo \\
GPU de consumo, 16\,GB & Sí & Holgado frente a la cota estimada \\
Partición MIG de 5\,GB & Al límite & Cercano a la cota de $4{,}2$\,GB; sensible al tamaño de sección \\
Una GPU de 48\,GB & Sí, referencia & Ninguno; margen amplio \\
Varias GPU, una sección & Sin beneficio & La profundidad del modelo es secuencial \\
Varias GPU, arreglo de secciones & Sí, recomendado & Número de secciones disponibles \\
Entrenamiento distribuido & No inmediato & Acoplamiento del lote en la asignación \\
\bottomrule
\end{tabular}
\end{table}

Dos filas requieren precisión. Añadir unidades no acelera una sección, porque
las etapas se ejecutan en orden estricto y el paralelismo aprovechable es
interno a cada capa, en parches y en cabezas, de modo que ya se agota dentro de
una sola unidad: la latencia de una sección es una constante del método. El
beneficio de añadir unidades se manifiesta únicamente en el tiempo total de
procesamiento del conjunto de secciones. Si $S$ es el número de secciones y $G$
el de unidades disponibles, la aceleración respecto de la ejecución secuencial es
\begin{equation}
A(G)=\min(G,\,S),
\label{eq:aceleracion}
\end{equation}
de modo que asignar más unidades que secciones no produce ninguna mejora
adicional. Este límite no proviene de una ineficiencia del método sino de su
granularidad: la sección es la unidad indivisible de trabajo. La consecuencia
práctica es que la partición de la catedral en secciones, adoptada inicialmente
por motivos de memoria, determina también el grado máximo de paralelismo
alcanzable.

\subsection*{D.5 Delimitación del aporte}
\label{sec:sintesis-analisis}

La Tabla~\ref{tab:atribucion} fija qué se modifica, qué se conserva y qué es
atribuible a este trabajo.

\begin{table}[H]
\centering
\caption{Delimitación del aporte a partir del análisis de los componentes.}
\label{tab:atribucion}
\begin{tabular}{@{}p{4.3cm}p{4.1cm}p{4.1cm}@{}}
\toprule
\textbf{Elemento} & \textbf{Origen} & \textbf{Estado en este trabajo} \\
\midrule
Codificador PTv3 y sus pesos & Sonata & Se emplea sin modificación \\
Paralelismo interno por parches y cabezas & PTv3 & Propiedad heredada, no aporte \\
Red troncal dispersa y cabezas de prototipos & HyperNCD & Se emplean sin modificación \\
Razonamiento por hipergrafo & HyperNCD & Se conserva íntegro \\
Descriptor geométrico de 6 dimensiones & HyperNCD & \textbf{Sustituido por $f_{\mathrm{ssl}}$} \\
Desacoplamiento del \emph{pipeline} y cacheo & Este trabajo & \textbf{Decisión propia} \\
Cota de aceleración, Ec.~\eqref{eq:aceleracion} & Este trabajo & \textbf{Formulación propia} \\
\bottomrule
\end{tabular}
\end{table}

El desacoplamiento del \emph{pipeline} merece una precisión: las dos ramas
exigen pilas de dependencias mutuamente excluyentes, de modo que ejecutar el
\emph{encoder} una sola vez por sección y persistir su salida resuelve esa
incompatibilidad, elimina la repetición del costo en cada época y habilita la
distribución sin sincronización de la Ec.~\eqref{eq:aceleracion}. Es una
decisión de ingeniería antes que de método, pero condiciona qué experimentos son
ejecutables con los recursos disponibles.
